\uuid{DWH2}
\titre{Pgcd}
\auteur{bodin ; croizat ; sacre}
\organisation{exo7}
\licence{CC BY-NC-SA 4.0}
\datecreate{2022-03-08}
\niveau{L1}
\module{Algèbre}
\chapitre{Arithmétique dans Z}
\sousChapitre{Pgcd, ppcm, algorithme d'Euclide}
\difficulte{1}
\theme{pgcd}
\exo7qcmid{hilisit.arithmetique.002}

\contenu{
\begin{question}
\qclassification{103.03}
    Quelles sont les valeurs qui correspondent à la division euclidienne $a=bq+r$ de $a$ par $b$ ?
    \begin{answers}
        \good{$a=48$, $b=7$, $q=6$, $r=6$}
        \good{$a=101$, $b=11$, $q=9$, $r=2$}
        \bad{$a=56$, $b=9$, $q=5$, $r=11$}
        \bad{$a=123$, $b=10$, $q=13$, $r=-7$}
    \end{answers}
    \begin{explanations}
    $48 = 7 \times 6 + 6$ \\
    $101 = 11 \times 9 + 2$ \\
    $56 = 9 \times 6 + 2$. Attention $56 = 9 \times 5 + 11$, mais on n'a pas $0 \le 11 \le 9-1$, donc cette écriture n’est pas la division euclidienne de $56$ par $9$.\\
    $123 = 10 \times 12 + 3$. Attention $123 = 10 \times 13 + (-7)$, mais on n'a pas $0 \le -7 \le 10-1$, donc cette écriture n’est pas la division euclidienne de $123$ par $10$.\\
    \end{explanations}
\end{question}
}
